\subsection{非可分几何下哪些结构会改变}
\label{sec:geometry-scope}

圆环不仅是便于描述的测试形状。其度量和边界条件不依赖角度，因此隐式算子按角模分块，非线性乘积则仍耦合各模态。设保留径向秩为 $d$，正角模数为 $K$，径向求积点数为 $q$，补零角向长度为 $N$，则独立模态的稠密分解准备量级为 $Kd^3$；每步包含量级 $Kd^2$ 的模态求解、$qdK$ 的重构与投影，以及 $qN\log N$ 的变换工作。这里略去固定通道数和独立均值块，只描述当前实现，不是一般RFM的复杂度定理。对于 $q\geq n$ 的稠密情形，具有 $n$ 个原始特征列的微分归一化还需要量级 $qn^2+n^3$ 的准备工作。

若映射边界随角度变化，度量系数及边界提升通常也依赖角度。与这些系数相乘会使隐式线性部分也耦合角模，独立模态方程及其各自的共同相位基不再构成全问题的精确分解。若直接对 $D=O(Kd)$ 个耦合自由度使用稠密求解，存储和分解量级将分别为 $O(D^2)$、$O(D^3)$。这是可以避免的稠密实现成本，不是必然下界：局部稀疏空间、无矩阵度量乘积、区域分解和迭代求解均可作为替代，但需要重新设计算子并验证，当前GPU测量没有检验这些方案。

可以在不夸大圆环变换作用的前提下明确条件性。设 $B$ 是同一离散空间上的对称正定参考隐式算子，$A$ 是映射域的对称正定算子。若几何及边界构造允许常数 $0<\alpha\leq\beta$ 满足
\begin{equation}
 \alpha x^\top Bx\leq x^\top Ax\leq\beta x^\top Bx
 \quad\text{对所有 }x,
 \label{eq:geometry-equivalence}
\end{equation}
则由Rayleigh商立即得到
\begin{equation}
 \sigma(B^{-1/2}AB^{-1/2})\subset[\alpha,\beta],\qquad
 \kappa_2(B^{-1/2}AB^{-1/2})\leq\beta/\alpha.
 \label{eq:geometry-condition}
\end{equation}
这只是有条件的预条件界。本文没有为任意物体证明与网格无关的 $\alpha,\beta$；映射畸变、窄缝及局部非均匀分辨率可能使其恶化。边界数据近似另行引入该谱界不能控制的一致性误差，却未必改变隐式矩阵。同样，圆环保留空间中的 $K_1=\nu I$ 是指定求积下的精确算术归一化恒等式，不会自动控制每个模态的质量矩阵、被截断方向的逼近误差或浮点重构误差。

使用单位分解的局部随机特征已用于复杂界面问题，包括线性Stokes算例\cite{r20}。这提供了可研究的扩展方向，但不能验证当前非稳态非线性求解器在这些几何上的表现。可信的扩展需要保持边界与散度约束，与确定性局部空间比较保留维数及求积工作，再重新完成速度、涡量和完整成本验收。因此，本文的紧凑空间与核性能结论限于指定的可分有限圆环。
